% ---------------------------------------------------------------------------
% Section: Phase 16 -- the code agent against modern harnesses
% ---------------------------------------------------------------------------
\section{Phase 16: The Code Agent Against Modern Harnesses}
\label{sec:phase16}

Phase~16 (2 November 2026) asks what the current agent harnesses and papers do that \Sys{}'s code agent does not, which of those
differences matter on four CPU cores with 1.7B and 4B models, and which can be closed and measured here. The research note is
\texttt{docs/research/harness\_gaps.md} and the decision is \texttt{docs/decisions/010-harness-gaps.md}.

\subsection{What modern harnesses do}

The sources fall into two groups. Four were read in full: the SWE-agent paper~\cite{yang2024sweagent}, the Agentless
paper~\cite{xia2024agentless}, the JetBrains study of observation masking~\cite{lindenbauer2025complexity}, and Anthropic's article on
context engineering~\cite{anthropic2025context}. The rest was read through search results (documentation pages, the DeepSeek-V3.2
report~\cite{deepseek2025v32}, Gemini CLI's documentation~\cite{geminicli2026docs}, the SWE-Search and SWE-Gym
papers~\cite{antoniades2025sweSearch,pan2024swegym}); a number from a vendor's own announcement is a vendor number.

\begin{itemize}
  \item \textbf{An edit-time guardrail} (SWE-agent's agent--computer interface). After every edit the file is checked, a failing edit
    is reverted, and the agent sees the error and the lines around it. On SWE-bench Lite with GPT-4 Turbo the resolve rate was 18.0\,\%
    with the check, 15.0\,\% without it and 10.3\,\% with no edit command at all~\cite{yang2024sweagent}.
  \item \textbf{Context management.} Compaction, dropping old tool results, notes kept outside the window and sub-agents that
    return a condensed result~\cite{anthropic2025context}; masking the observations older than ten turns halves the cost with the
    solve rate of LLM summarization across five model configurations~\cite{lindenbauer2025complexity}.
  \item \textbf{Selection by a verifier.} Majority voting resolved 25.67\,\% of SWE-bench Lite, regression tests added 27.00\,\% and
    reproduction tests 32.00\,\%~\cite{xia2024agentless}; trained critics rerank trajectories~\cite{pan2024swegym}; search backtracks to a
    better node~\cite{antoniades2025sweSearch}.
  \item \textbf{Checkpoints and project instructions.} Gemini CLI snapshots the project before file-modifying actions and reads a
    hierarchy of \texttt{GEMINI.md} files; Codex and Claude Code read \texttt{AGENTS.md} and \texttt{CLAUDE.md}.
  \item \textbf{Stall detection.} Loop guards key on the tool name, its arguments and its result; repetition warns, repetition with
    stagnant results stops. This is engineering practice; no study was found.
  \item \textbf{Reasoning persistence.} DeepSeek-V3.2 keeps its reasoning across the tool calls of one turn and drops it when a new user
    message arrives~\cite{deepseek2025v32}.
\end{itemize}

\subsection{The gaps}

\Sys{} already had an append-only transcript that keeps the server's prompt cache valid, compaction written by the harness, a patch
grammar, Codex's patch parser, an \emph{exact} verifier after every change (a stronger selection signal than a critic model or
generated tests), the cascade, retries and a wall-clock cap. Four gaps were real and are closed. A file that did not compile went to the
verifier like any other, costing a sandbox run and a whole step, and the same \texttt{IndentationError} came back step after step. A step
that made the tests do worse was kept as the base of the next one. Only an identical \emph{workspace} was noticed as a loop, so the same
failure from different code, or a reply with no action, repeated until the step budget (one run spent the whole 900\,s cap that way).
And the agent took the task text only.

Others are named and not closed, with the reason. Self-generated reproduction tests need a bug that fails them on the unpatched code; a new
function has no such oracle, and a test written by a 1.7B model would raise the \emph{label} without raising the trust. Reasoning persistence
is not adopted now: the Qwen3-Thinking template (read from the model file) keeps an assistant turn's reasoning only after the last user
query, and the harness sends observations as ordinary user messages, so each step is a new query; keeping the reasoning would mean sending
tool responses in the template's form and about 25\,s of prefill per kept 500-token span at the measured 20 tokens/s, and is the next
experiment. Sub-agents, plan mode and hooks have no counterpart at the scale of one function in a flat workspace, where the cascade is the
sub-agent idea and the sandbox is a hook that cannot be skipped. Observation masking changes nothing inside a ten-turn window, and the runs
are at most twelve steps.

\subsection{What was built}

\begin{enumerate}
  \item \textbf{A compile check before the verifier} (\texttt{harness.lint}). After a patch or a whole-file write every edited Python file is
    compiled in the sandbox (never executed). A file that does not compile is \emph{not applied}: the files stay, and the model is shown the
    error and the five lines around it with the offending one marked. No verifier run is spent, and a cascade candidate that does not
    compile is not judged. A failure of the sandbox itself passes the edit on to the verifier.
  \item \textbf{Rollback to the best version} (\texttt{harness.rollback}). A change judged below the best version so far is dropped, the
    files are restored to that version, and the model is told what was tried; the failed attempt stays in the transcript.
  \item \textbf{A stall guard} (\texttt{harness.stuck\_after}, default 4). A failed step's signature is the hidden tests' evidence, or the
    action and its first line of feedback when the files were not judged. The second consecutive identical failure carries a nudge to change
    approach; the $k$-th ends the run with status \texttt{stuck} and the best files.
  \item \textbf{Caller guidelines} (\texttt{guidance}, 2,000 characters), shown with the task and again after a compaction.
\end{enumerate}

Each is a request field and a configuration key, on by default, and reported in the reply (\texttt{lint\_rejections}, \texttt{rollbacks},
the step's \texttt{rolled\_back}, the status \texttt{stuck}). Five unit tests drive the loop with a scripted host, and an integration test
runs the compile check through the real sandbox.

\subsection{Measurements}

All runs were on the shared host (the owner's other workloads at load 5--15) with Smiley Core, over the same ten tasks (eight repairs and writes of one
function with hidden tests), the arms interleaved and their order alternated per task: guards off (\texttt{lint}, \texttt{rollback},
\texttt{stuck\_after}\,=\,0) against the production defaults.

\begin{center}
\begin{tabular}{llrrr}
\hline
Set & Arm & Verified & Model calls & Wall time \\
\hline
System 2, 6 steps & off & 10/10 & 17 & 747\,s \\
 & on & 10/10 & 17 & 772\,s \\
System 1, 8 steps, round 1 & off & 4/10 & 58 & 307\,s \\
 & on & 4/10 & 58 & 297\,s \\
System 1, 8 steps, round 2 & off & 5/10 & 47 & 350\,s \\
 & on & 5/10 & 47 & 347\,s \\
\hline
\end{tabular}
\end{center}

The guards cost nothing: the verified counts and model calls are identical, the decoded tokens are equal except for one System 1 task whose
sampling differed (2,418 against 2,362 in the first round), and the wall times agree within the host's noise. The compile check costs 37\,ms per edit (35\,ms of it Python's start-up) against steps of 5 to 230\,s. They also changed no
outcome here: neither model produced a file that does not compile, and rollback restored nothing, so the benefit lies in the failure modes
the research names, which this set rarely produces. Those are pinned by tests and were seen once in practice, as a Smiley Core run that
spent its whole 900\,s cap on repeated \texttt{IndentationError}s. The A/B did expose a defect in the first version of the stall guard:
System 1 asked six times in a row to run a file that does not exist, and the guard skipped \texttt{run} steps. It now counts them, hashes
only the stable lines of the feedback, and the message lists the files that exist. The same ten System 1 tasks verified four times in the
first round and five in the second in both arms, so a one-task difference is sampling noise.

\subsection{Summary}

The research found one capability the code agent lacked that matters at this scale, the edit-time guardrail, and three smaller ones
(checkpoint and restore, stall detection, project instructions); all four are built, on by default, tested, and measured to cost nothing.
Self-generated tests, reasoning persistence and sub-agents are named as gaps with the reason they were not closed. The verifier-based
selection that other systems approximate with critic models is the one place \Sys{} is ahead: its verifier is exact.
